
\section{Question: what does torque actually tell us about a flux vector?}
The familiar dq torque equation is usually used in the forward direction: flux and current predict torque. For flux-map validation the inverse question is more revealing. Given current and an independent torque observation, which part of the two-dimensional flux vector has actually been measured?

The answer is a projection, not a complete reconstruction. This paper develops that geometry, derives the corresponding solution line and minimum-change representative, and then asks whether magnetic reciprocity removes the missing direction. It does not: smooth radial potential terms remain invisible to torque. The synthetic and conditional real-data studies are organized around this observability statement. Co-energy fitting and general measurement-error identification are deliberately kept in separate papers.
